
\documentclass{article}
\usepackage{colortbl}
\usepackage{makecell}
\usepackage{multirow}
\usepackage{supertabular}

\begin{document}

\newcounter{utterance}

\centering \large Interaction Transcript for game `bbh', experiment `formal\_fallacies', episode 56 with BSU{-}SLIM\_\_Qwen3.5{-}9B{-}Base\_\_playpen{-}prm{-}9b{-}best\_of\_n.
\vspace{24pt}

{ \footnotesize  \setcounter{utterance}{1}
\setlength{\tabcolsep}{0pt}
\begin{supertabular}{c@{$\;$}|p{.15\linewidth}@{}p{.15\linewidth}p{.15\linewidth}p{.15\linewidth}p{.15\linewidth}p{.15\linewidth}}
    \# & \multicolumn{2}{c}{Player} && \multicolumn{2}{c}{Game Master} \\
    \hline

    \theutterance \stepcounter{utterance}  
    & & & \multicolumn{4}{p{0.6\linewidth}}{
        \cellcolor[rgb]{0.9,0.9,0.9}{
            \makecell[{{p{\linewidth}}}]{
                \texttt{\tiny{[P1$\langle$GM]}}
                \texttt{Distinguish deductively valid arguments from formal fallacies.} \\
\\ 
\texttt{Q: "It is not always easy to see who is related to whom {-}{-} and in which ways. The following argument pertains to this question: To begin with, Lesley is a close friend of Fernando. Moreover, being a close friend of Fernando or a schoolmate of Lowell is sufficient for being a great{-}grandfather of Leroy. It follows that Lesley is a great{-}grandfather of Leroy."} \\
\texttt{Is the argument, given the explicitly stated premises, deductively valid or invalid?} \\
\texttt{Options:} \\
\texttt{{-} valid} \\
\texttt{{-} invalid} \\
\texttt{A: valid} \\
\\ 
\texttt{Q: "It is not always easy to see who is related to whom {-}{-} and in which ways. The following argument pertains to this question: Whoever is not a great{-}grandfather of Clyde is a stepbrother of Brian. Being an ancestor of Dana is sufficient for not being a great{-}grandfather of Clyde. We may conclude: Everyone who is an ancestor of Dana is a stepbrother of Brian, too."} \\
\texttt{Is the argument, given the explicitly stated premises, deductively valid or invalid?} \\
\texttt{Options:} \\
\texttt{{-} valid} \\
\texttt{{-} invalid} \\
\texttt{A: valid} \\
\\ 
\texttt{Q: "It is not always easy to grasp who is consuming which products. The following argument pertains to this question: Every infrequent user of Paul Mitchell shampoo is either a rare consumer of Nioxin shampoo or a loyal buyer of Caress soap, or both. No regular consumer of Lush soap is a rare consumer of Nioxin shampoo and, in the same time, a loyal buyer of Caress soap. It follows that whoever is an infrequent user of Paul Mitchell shampoo is not a regular consumer of Lush soap."} \\
\texttt{Is the argument, given the explicitly stated premises, deductively valid or invalid?} \\
\texttt{Options:} \\
\texttt{{-} valid} \\
\texttt{{-} invalid} \\
\texttt{A: invalid} \\
\\ 
\texttt{Q: "Here comes a perfectly valid argument: First, everyone who is a half{-}brother of Ronny or a classmate of Sal is not a son of Ronald. Second, it is false that Santo is a son of Ronald. So, necessarily, Santo is a half{-}brother of Ronny or a classmate of Sal."} \\
\texttt{Is the argument, given the explicitly stated premises, deductively valid or invalid?} \\
\texttt{Options:} \\
\texttt{{-} valid } \\
\texttt{{-} invalid} \\
\texttt{A:} \\
            }
        }
    }
     \\ \\

    \theutterance \stepcounter{utterance}  
    & \multicolumn{4}{p{0.6\linewidth}}{
        \cellcolor[rgb]{0.561,0.737,0.561}{
            \makecell[{{p{\linewidth}}}]{
                \texttt{\tiny{[P1$\rangle$GM]}}
                \texttt{This argument is **invalid**.} \\
\\ 
\texttt{Here is the step{-}by{-}step logical breakdown:} \\
\\ 
\texttt{1.  **Analyze the Premises:**} \\
\texttt{    *   **Premise 1:** \$(H \textbackslash{}lor C) \textbackslash{}rightarrow \textbackslash{}neg S\$} \\
\texttt{        *   (If someone is a half{-}brother of Ronny OR a classmate of Sal, THEN they are NOT a son of Ronald).} \\
\texttt{        *   This establishes a sufficient condition: Being a half{-}brother or classmate guarantees that the person is *not* a son.} \\
\texttt{    *   **Premise 2:** \$\textbackslash{}neg S\$ is true for Santo.} \\
\texttt{        *   (Santo is NOT a son of Ronald).} \\
\\ 
\texttt{2.  **Analyze the Conclusion:**} \\
\texttt{    *   **Conclusion:** \$H \textbackslash{}lor C\$ is true for Santo.} \\
\texttt{        *   (Santo is a half{-}brother of Ronny OR a classmate of Sal).} \\
\\ 
\texttt{3.  **Identify the Logical Error:**} \\
\texttt{    *   The argument attempts to use the fallacy of the **fallacy of the inverse** (or confusing the inverse with the converse).} \\
\texttt{    *   From Premise 1, we know that if the antecedent (\$H \textbackslash{}lor C\$) is true, then the negation of the consequent (\$\textbackslash{}neg S\$) must be true.} \\
\texttt{    *   However, simply knowing that the *consequent* (\$\textbackslash{}neg S\$) is true does **not** prove that the *antecedent* (\$H \textbackslash{}lor C\$) is true.} \\
\texttt{    *   **Counter{-}example:** Santo might be neither a half{-}brother nor a classmate. In that case, Santo is still not a son of Ronald (perhaps Santo is an estranged brother, or an adopted child, or just a random person who happened not to be a son). The premises do not rule out other reasons why someone might not be a son of Ronald.} \\
\\ 
\texttt{4.  **Formal Logic Check:**} \\
\texttt{    *   Let \$P\$ = is half{-}brother or classmate.} \\
\texttt{    *   Let \$Q\$ = is son.} \\
\texttt{    *   Premise: \$P \textbackslash{}rightarrow \textbackslash{}neg Q\$} \\
\texttt{    *   Premise: \$\textbackslash{}neg Q\$ (for Santo)} \\
\texttt{    *   Conclusion: \$P\$} \\
\texttt{    *   The structure is \$P \textbackslash{}rightarrow \textbackslash{}neg Q\$ and \$\textbackslash{}neg Q\$, therefore \$P\$. This is a classic invalid form (\$Denying the Antecedent\$ if phrased differently, or affirming the Consequent in the contrapositive sense). Specifically, it is invalid because the truth of the consequent does not guarantee the truth of the antecedent.} \\
\\ 
\texttt{**Answer:**} \\
\texttt{invalid} \\
            }
        }
    }
    & & \\ \\

    \theutterance \stepcounter{utterance}  
    & & & \multicolumn{2}{p{0.3\linewidth}}{
        \cellcolor[rgb]{0.9,0.9,0.9}{
            \makecell[{{p{\linewidth}}}]{
                \texttt{\tiny{[GM$|$GM]}}
                \texttt{Target: invalid} \\
            }
        }
    }
    & & \\ \\

    \theutterance \stepcounter{utterance}  
    & & & \multicolumn{2}{p{0.3\linewidth}}{
        \cellcolor[rgb]{0.9,0.9,0.9}{
            \makecell[{{p{\linewidth}}}]{
                \texttt{\tiny{[GM$|$GM]}}
                \texttt{game\_result = LOSE} \\
            }
        }
    }
    & & \\ \\

\end{supertabular}
}

\end{document}
